% AUTORIDAD DE FUENTE: Exposé IX, §§5.7--5.10, índices de escaneo
% 392--397, folios 381--386. No hay diferencias materiales D/C en este
% componente; los subíndices o mecanografiados se restituyen como 0.

%%%%%%%%%%  scan idx 392  |  folio 381  %%%%%%%%%%

\par\medskip
\noindent
\underline{Proposición} 5.7. \underline{Supongamos que $A^0$ es
$\ell$-divisible y que $S$ es henseliano. Sea $F_K$ un haz
$\ell$-ádico constructible sobre $K$ que sea ``no ramificado'', es
decir, que se prolongue a un haz $\ell$-ádico constante torcido $F$
sobre $S$, y sea $u:F_K\longrightarrow T_\ell(A_K)$ un homomorfismo de
progrupos. Entonces $u$ se factoriza por un homomorfismo de $F$ a
$T_\ell(A_K)^f$.}

En efecto, escribamos
\begin{equation}\tag{5.7.1}
 u=(u(\nu))_{\nu\geqslant0},\qquad
 u(\nu):F_K(\nu)\longrightarrow T_\ell(A_K)(\nu)
 ={}_{\ell^\nu}A_K.
\end{equation}
Por la propiedad universal de $A$, como $F(\nu)$ es étale y, por
tanto, liso sobre $S$, el morfismo $u(\nu)$ procede de un homomorfismo
\[
 v(\nu):F(\nu)\longrightarrow{}_{\ell^\nu}A.
\]
Por definición de la parte fija, este homomorfismo se factoriza como
\[
 v(\nu):F(\nu)\longrightarrow({}_{\ell^\nu}A)^f.
\]
Como la imagen del morfismo inducido sobre las fibras especiales está
contenida en la parte infinitamente $\ell$-divisible de $A_0$, se sigue
que $v(\nu)$ se factoriza de hecho por
$({}_{\ell^\nu}A^0)^f$. Pasando al límite sobre $\nu$ se obtiene un
homomorfismo
\[
 v:F\longrightarrow T_\ell(A^0)^f,
\]
y el morfismo inducido sobre las fibras genéricas da manifiestamente la
factorización buscada.

\medskip
\noindent
\underline{Observaciones} 5.7.2. Si $\ell\neq p$, la propiedad 5.7
caracteriza de manera única el subprogrupo $T_\ell(A_K)^f$ de
$T_\ell(A_K)$ como su mayor subhaz $\ell$-ádico no ramificado. Si
$\ell=p$, el enunciado de 5.7 deja de ser muy satisfactorio. Daremos
una propiedad considerablemente más fuerte, que podrá servir como
caracterización intrínseca de la parte fija de $T_p(A_K)$ —en términos
de la sola estructura de progrupo de Barsotti--Tate de $T_p(A_K)$—, al
menos en características desiguales.

%%%%%%%%%%  scan idx 393  |  folio 382  %%%%%%%%%%

\medskip
\noindent
\underline{Teorema} 5.8 (M. Raynaud). \underline{Sea $A_K$ un esquema
abeliano con reducción semiestable sobre el cuerpo de funciones $K$ de
un trait henseliano $S$. Sea $\ell$ un número primo (el caso
interesante es $\ell=p$), y consideremos el progrupo de
Barsotti--Tate pro-$\ell$ $X_\eta=T_\ell(A_K)$ y su ``parte fija''
$X_\eta^f$, que es la fibra genérica del progrupo de Barsotti--Tate
pro-$\ell$ $X^f=T_\ell(A^0)^f$. Si $\ell=p$, supongamos que
$\operatorname{car}K=0$ (para poder aplicar sobre $S$ el teorema de
Tate [33, Teor.\ 4]). Sea $F$ un progrupo de Barsotti--Tate pro-$\ell$
sobre $S$. Entonces todo homomorfismo
$F_\eta\longrightarrow X_\eta$ se factoriza por $X_\eta^f$.}

Pasando a los grupos ind-$\ell$ asociados con los progrupos de
Barsotti--Tate pro-$\ell$ considerados, el enunciado precedente se
vuelve equivalente, por el teorema citado de Tate, al siguiente.

\medskip
\noindent
\underline{Corolario} 5.9. \underline{Si $F$ es un grupo de
Barsotti--Tate ind-$\ell$, entonces todo homomorfismo
$u_\eta:F_\eta\longrightarrow A_\eta$ se prolonga de manera única a un
homomorfismo $u:F\longrightarrow A$.}

En otras palabras, la propiedad néroniana característica (1.1.2) del
modelo de Néron sigue siendo válida para homomorfismos procedentes de
ind-esquemas en grupos de tipo Barsotti--Tate, sin hipótesis de lisura
sobre ellos. Solo el caso de grupos ind-étale se deduciría directamente
de la definición.

%%%%%%%%%%  scan idx 394  |  folio 383  %%%%%%%%%%

\medskip
\noindent
\underline{Observación} 5.9.1. En realidad, M. Raynaud demuestra 5.9
sin suponer que $A_K$ tenga reducción semiestable, e incluso para todo
esquema en grupos néroniano $A$ sobre $S$ (cuya fibra genérica no tiene
por qué ser un esquema abeliano). Naturalmente, la hipótesis de que $S$
sea henseliano tampoco es necesaria en la formulación de 5.9. La
demostración de Raynaud en este caso más general lo reduce al caso
particular considerado en 5.8 mediante una propiedad inédita de la
construcción de los modelos de Néron. Aquí daremos únicamente la
demostración de Raynaud de 5.9.

\medskip
\noindent
\underline{Demostración de 5.9.} Se sabe [33] que $F$ es una extensión
de un grupo de Barsotti--Tate ind-étale $F'$ por un grupo de
Barsotti--Tate $F''$ de fibra especial conexa, de modo que para todo
entero $n$ existe una sucesión exacta
\[
 0\longrightarrow F''(n)\longrightarrow F(n)
 \longrightarrow F'(n)\longrightarrow0,
\]
con $F'(n)$ étale y $F''(n)_0$ puramente infinitesimal. El homomorfismo
dado $u_\eta$ está definido por una sucesión de homomorfismos
\[
 u_\eta(n):F(n)_\eta\longrightarrow A_\eta,
\]
y basta prolongarlos por separado a homomorfismos
$u(n):F(n)\longrightarrow A$. El Lema 5.9.2 reduce entonces el problema
a demostrar que el homomorfismo
$u''_\eta:F''_\eta\longrightarrow A_\eta$ inducido por $u_\eta$ se
prolonga a $u'':F''\longrightarrow A$. Por consiguiente,
\underline{podemos suponer que $F$ tiene fibra especial conexa}.

Podemos considerar $u_\eta$ como un homomorfismo de $F_\eta$ al grupo
de Barsotti--Tate ind-$\ell$ $Y_\eta$ asociado a
$X_\eta=T_\ell(A_K)$, y considerar después el compuesto
\begin{equation}\tag{*}
 v_\eta:F_\eta\longrightarrow Y_\eta
 \longrightarrow Y_\eta/Y_\eta^f.
\end{equation}

%%%%%%%%%%  scan idx 395  |  folio 384  %%%%%%%%%%

Por 5.6, este último grupo se prolonga a un grupo de Barsotti--Tate
ind-étale $Z$ sobre $S$, a saber,
\[
 \underline M\otimes_{\mathbb Z}(\mathbb Q_\ell/\mathbb Z_\ell).
\]
Por el teorema de Tate citado en 5.8, el homomorfismo (*) se prolonga a
un homomorfismo
\[
 v:F\longrightarrow Z.
\]
Como $F$ tiene fibra especial conexa y $Z$ es ind-étale, este morfismo
es nulo. Por tanto, $v_\eta$ es también nulo y $u_\eta$ se factoriza por
$Y_\eta^f$, como se quería.

Queda por demostrar el lema siguiente.

\medskip
\noindent
\underline{Lema} 5.9.2. Sea $S$ un trait, sea $A$ un esquema en grupos
néroniano sobre $S$ (1.1.2), sea
\[
 1\longrightarrow G'\longrightarrow G\longrightarrow G''
 \longrightarrow1
\]
una extensión de esquemas en grupos sobre $S$, con $G''$ liso, sea
$u:G'\longrightarrow A$ un homomorfismo de grupos, y sea
$v_\eta:G_\eta\longrightarrow A_\eta$ un homomorfismo de grupos que
prolonga $u_\eta$. Supongamos que la extensión de $G''$ por $A$
obtenida de la extensión $G$ por pushout mediante $u$ es representable
\footnote{Resultados inéditos de M. Raynaud implican, además, que esta
condición siempre se cumple.} —como ocurre, por ejemplo, si
$G\longrightarrow G''$ admite localmente una sección para la topología
``finita localmente libre'' (véase SGA 3, IV 6.3) y si $A$ es
cuasiproyectivo sobre $S$. Entonces existe un único homomorfismo de
grupos $v:G\longrightarrow A$ que prolonga a la vez $v_\eta$ y $u$.

%%%%%%%%%%  scan idx 396  |  folio 385  %%%%%%%%%%

En efecto, la última hipótesis nos reduce al caso en que $G'=A$, $u$
es la identidad y $v_\eta:G_\eta\longrightarrow A_\eta$ es una
proyección que se quiere prolongar a una proyección
$v:G\longrightarrow A$. Pero $G$ es liso, por ser extensión de dos
esquemas en grupos lisos, y la conclusión se sigue de la definición de
la propiedad néroniana de $A$.

\medskip
\noindent
\underline{Corolario} 5.10. \underline{Sea $A_K$ un esquema abeliano
sobre el cuerpo de funciones} $K$ \underline{de un trait $S$, y sea $\ell$ un
número primo. Si $\ell=p$, supongamos que $\operatorname{car}K=0$, de
modo que pueda aplicarse el teorema de Tate ya invocado en 5.8. Entonces
$A_K$ tiene buena reducción sobre $S$ si y solo si el progrupo de
Barsotti--Tate pro-$\ell$ $T_\ell(A_K)$ ``tiene buena reducción sobre
$S$'', es decir, si se prolonga a un grupo de Barsotti--Tate sobre $S$.}

(Este es el resultado conjeturado en [8, 4.9].)

La necesidad es clara; demostremos la suficiencia. Obsérvese que el caso
$\ell\neq p$ ya fue tratado en 2.2.9. Podemos suponer que $S$ es
henseliano (2.2.9). Supongamos primero que $A_K$ tiene reducción
semiestable sobre $S$. Aplicando 5.8 con $F=X$ se obtiene
\[
 T_\ell(A_K)=T_\ell(A_K)^f.
\]
Esto significa también que el toro maximal de $A_0$ es cero, es decir,
que $A_0^0$ es un esquema abeliano; la conclusión se sigue entonces de
2.2.9 (iii${}^0$).

%%%%%%%%%%  scan idx 397  |  folio 386  %%%%%%%%%%

Tratemos ahora el caso general. Por 2.2.9 podemos suponer que $S$ es
estrictamente local. Sea $K'/K$ una extensión de Galois finita tal que
$A_{K'}$ tenga reducción semiestable sobre la normalización $S'$ de $S$
en $K'$ (3.6). Si $A^*$ denota el modelo de Néron correspondiente,
$G=\operatorname{Gal}(K'/K)$ actúa sobre la situación $(A^*,S')$ y,
por tanto, sobre $A_0^{*0}$. Por 4.1 basta demostrar que $G$ actúa
trivialmente sobre este último; ello implica, en efecto, que $A_K$ tiene
reducción semiestable y nos reduce al caso ya tratado.

Sea $X$ un grupo de Barsotti--Tate sobre $S$ que prolonga
$X_\eta=T_\ell(A_K)$. Por el teorema de unicidad de Tate ya invocado,
existe un único isomorfismo
\[
 X\times_S S'\simeq T_\ell(A^{*0})
\]
que induce la identidad sobre las fibras genéricas. Por transporte de
estructura, este isomorfismo es compatible con las acciones de $G$.
Ahora bien, $G$ actúa trivialmente sobre la fibra especial de
$X\times_S S'$ y, por tanto, sobre $T_\ell(A_0^{*0})$. Se sigue
inmediatamente que $G$ actúa trivialmente sobre $A_0^{*0}$, lo que
concluye la demostración.
